\newpage
% \color{white}
% \pagecolor{black}

\ifdefined\useChapters
\chapter{O vilarejo encantado}
\else
\chapter{}
\fi

Terça-feira de manhã. O sol já alto lançava seus raios sobre as ruas movimentadas da cidade, onde o som dos carros e o murmúrio das conversas enchiam o ar. Caminhava pelas calçadas, observando as vitrines das lojas que se preparavam para mais um dia comum. No entanto, para mim, nada parecia comum. Era como se a cidade tivesse esquecido o corpo inerte da mulher na praça central, mas eu não conseguia tirar aquela imagem da cabeça. Seus olhos sem vida e a expressão de terror me perseguiam, um aviso silencioso para tomar cuidado. Eu sabia que havia alguém atrás de mim também, tal como havia estado atrás dela.

Não sei do que são capazes ou mesmo quem são, mas uma coisa é certa: quando me encontrarem, porque sim, eu sei que vão, preciso estar preparado para fugir ou revidar. Minha mente trabalhava freneticamente, tentando encontrar uma solução. Talvez eu pudesse criar uma ilusão, esconder-me atrás de um véu de mentiras visuais, tal como provavelmente ela tentou fazer. Mas se ela, com suas habilidades, não teve sucesso, por que eu teria? A realidade era fria e dura. Quem quer que fosse, sabia muito bem o que estava fazendo, e eu sentia um arrepio gelado subir pela espinha ao pensar no que poderiam fazer comigo.

No caminho para a escola, encontro meus amigos, mas não consigo tirar os pensamentos sombrios da cabeça. Enquanto caminhávamos juntos, rindo e conversando sobre trivialidades, minha mente continuava a traçar planos de fuga. Talvez eu pudesse me esconder em grandes multidões, tal como nos filmes de espionagem. Parecia uma boa ideia; eles não iriam querer ser descobertos em público. Mas e se isso não fosse suficiente? E se eles me encontrassem mesmo assim?

— E aí, cara, você parece meio avoado hoje. — disse Jonas, o garoto problema da escola, enquanto escondia alguma coisa dentro da mochila com um movimento furtivo. Não sei por que, mas ele sempre foi um ímã para encrenca. Se houvesse algo grande acontecendo na escola, podia ter certeza de que Jonas estava envolvido, e na maioria das vezes, ele era o responsável. Seus olhos, sempre brilhando com uma mistura de malícia e entusiasmo, agora me observavam com curiosidade.

— Ah, problemas em casa. Sabe como é, minha irmã sempre em cima. — respondi, tentando soar casual, embora minha mente estivesse a mil.

— Aaaaah, sua irmã. Sua irmã é da hora! — disse Jonas com um sorriso malicioso.

— Que isso, mano? Deixa minha irmã em paz. — retruquei, sentindo o desconforto crescer.

— Meio difícil, viu — ele respondeu, rindo de canto de boca, seus olhos fixos nos meus de forma intimidadora. Havia algo nele que sempre me incomodava, uma sensação de que ele sabia mais do que deixava transparecer. Eu nunca gostei desse cara, sempre foi um grande babaca.

— Caramba, você viu o jogo de ontem? Aquele cara corre muito, né? — falei rapidamente, tentando puxar um assunto aleatório com a primeira pessoa que passou perto de mim. A tensão com Jonas era palpável, e eu precisava de uma distração.

— Que jogo, mano, você tá maluco? E não fala comigo, seu perdedor. Não quero que ninguém me veja com um lunático igual você. — retrucou o garoto, me lançando um olhar de desprezo.

Suspirei aliviado quando percebi que minha aula era no segundo andar, bem longe desse merda. Subi as escadas rapidamente, tentando afastar o desconforto.

Enquanto subia as escadas, notei um silêncio repentino e uma sensação como se o ar tivesse ficado mais denso. Era uma sensação familiar, uma premonição de que algo grande estava para acontecer. Parei por um momento, os pés hesitando no degrau, e olhei para baixo. Do corredor do segundo andar, consegui ver o pátio, onde uma movimentação estranha estava ocorrendo. Muitas pessoas estavam juntas, visivelmente espantadas e murmurando entre si.

Lá embaixo, no meio da confusão, estavam quatro pessoas vestidas de preto, três homens e uma mulher de cabelo azul curto. Eles tinham rendido Jonas e mais alguns alunos, e os homens começaram a gritar, suas vozes ecoando pelo pátio.

— Nós sabemos que você está aí! — gritou um sujeito calvo, um pouco mais velho que os outros.

— Não adianta se esconder! — continuou outro, a voz carregada de ameaça.

— Saia daí ou você vai ser o responsável por tudo isso! — acrescentou o terceiro homem, enquanto a garota de cabelo azul observava em silêncio, seus olhos varrendo o pátio com uma confiança inquietante.

Lá de baixo, a figura de cabelos azuis me olhou diretamente nos olhos. Era como se ela pudesse ver através de mim, lendo meus pensamentos e medos. Não disse uma palavra, mas seu olhar dizia tudo: ela já sabia que eu estava ali, esperando apenas o momento certo para agir. Sua confiança era perturbadora, como se tivesse certeza absoluta de que aquilo tudo acabaria bem rápido.

Desço correndo as escadas, o coração batendo descompassado. Apesar de todas as desavenças, Jonas não merecia ser estrangulado por minha causa. Nenhum dos alunos ali merecia pagar pelo que eu tinha feito ou pelo que eu era. Eu precisava fazer algo.

— Então é você que vem causando problemas por aqui? — perguntou o homem calvo, com um sorriso sinistro.

— É você que vem usando seus poderes. — disse outro, aproximando-se de mim.

— Só a gente tem direito. — acrescentou a garota de cabelo azul, sua voz fria e calculista.

Enquanto eles falavam, soltaram todos os reféns e, logo que todos saíram correndo do pátio, das mãos da menina de cabelo azul brotou uma enorme labareda que ela, mais rápido do que eu podia imaginar, lançou na minha direção. A chama não chegou a me atingir, mas pude sentir o calor, o que me espantou muito. Eu estava certo de que ela também estava criando uma ilusão, mas tive certeza de que não era, porque, como eu tinha consciência que talvez fosse uma ilusão, achei que ela não poderia mantê-la e, principalmente, porque quase me queimou.

— Tá calorzinho, tá? — me provocou a menina de cabelos azuis enquanto lançava mais algumas chamas na minha direção.

— Eu falei que ele era um bundão, não sei por que o velho dá tanto crédito — disse a menina enquanto o senhor mais velho vinha rindo na minha direção.

— É sempre assim, alguns recebem crédito enquanto a gente faz todo o trabalho sujo — disse o senhor enquanto vinha na minha direção.

Tentei fugir em direção ao portão, mas um dos três me bloqueou como se conseguissem me segurar, mesmo estando a alguns metros de mim, me fazendo cair para trás. Eles podiam fazer coisas que influenciavam os outros de fato e não só na minha mente.

Nunca senti tanto desespero na minha vida. No meu corpo, era como se eu estivesse descendo uma montanha-russa e não conseguisse saber se ia cair ou não.

Enquanto eu estava no chão, não pude deixar de notar que acima de mim estava lá ele, flutuando, meu antigo conhecido, e apesar de trazer no rosto uma feição de arrependimento, jogava pedras em mim. Consegui me livrar de algumas, mas elas eram muitas.

Meu Deus, existem muitas habilidades que eu não fazia ideia. Até então, achava que só havia as que eu tinha sentido e aprendido.

Enquanto eu tentava me defender, percebi que também conseguia evitar que as pedras chegassem em mim. Imaginei como se minhas mãos estivessem muito mais distantes do meu corpo do que meu braço permitia, e assim conseguia tocar coisas de longe. Consegui desviar as pedras que o garoto voador me lançava e, como no centro do pátio havia um pequeno chafariz, consegui movimentar uma quantidade de água, o máximo que pude, e joguei na menina de fogo, o que fez com que ela perdesse um pouco a concentração e o fogo se dissipasse.

— Que merda é essa, Larissa? Você me disse que ele era um ilusionista! — gritou o rapaz que tentava me segurar à distância.

— É melhor a gente acabar logo com isso, talvez ele seja mais do que imaginávamos.

Não vou conseguir aguentar lutar com eles por muito tempo, e aqui na escola alguém pode se ferir. Além disso, vendo todo esse uso de habilidades, alguém pode muito bem entender que também é capaz e compreender como usar algumas daquelas habilidades, e isso, como já aprendi, pode não ser uma boa ideia.

Como consegui despistar momentaneamente os três, vou sair correndo pelo portão e tentar me esconder na floresta que fica ao lado da escola. Estranho que, apesar de todos terem várias habilidades especiais, aquele senhor mais velho parecia não ter nada de especial, mas também não pretendo ficar pra descobrir.

Como sempre matei aula pra me esconder naquela floresta, tenho uma vantagem em cima deles.

— Vai por cima e pega ele, Ítalo.

— Você não manda em mim, quem foi que colocou você no comando?

— Sabe muito bem que se não fosse por mim estaria fazendo um chazinho praquele velho e daqui a pouco não seríamos nada de novo.

Saio correndo pela floresta sem olhar muito para onde vou, atrás de mim só escuto os barulhos de galhos se quebrando e chamas no ar.

— Não adianta você correr, seu medroso. Eu vou queimar essa floresta inteira até te achar.

Corro, corro. Como nunca havia corrido antes. Não consigo me concentrar o suficiente para poder criar alguma ilusão ou ficar invisível, estou com muito medo, muito agitado.

— Calma! Respira fundo e se concentra!

— Eu vou te pegar! — gritou Larissa, enfurecida.

— Não tem para onde ir — disse Ítalo, dando alguns rasantes entre as copas das árvores.

— Calma! Respira fundo.

Eu sei que é um momento de desespero, mas começo a pensar que essas habilidades trazem mais problemas do que soluções. Tantas pessoas se machucando. Tanto caos, e a troco de quê? De suprir as minhas necessidades?

Por um curto espaço de tempo, consigo me concentrar e acho que estou invisível. Eu nunca sei, porque mesmo invisível continuo vendo meu corpo.

Eu acho que sim, quando viro para trás, pela expressão em seus rostos, eles parecem ter me perdido de vista.

— Aaaaaaaah! O que é isso?

— Que droga! — falo enquanto tropeço em algum lugar e saio girando ladeira abaixo. Não me lembro de ter uma ladeira aqui. Devo estar bem mais longe do que já fui.

— Ai, minha cabeça!

— Ai, ai, ai.

Giro por um bom tempo, já nem sei quantas vezes bati minhas costas, minha cabeça.

— Ufa, acho que finalmente cheguei a algum lugar, mas não consigo mexer meu corpo!

Minha vista está meio embaçada, acho que foram de tantas pancadas. Estou um pouco tonto também, acho que vou desmaiar.

— Pera aí!

— Aaah! O que é isso, é um vilare....

Desmaiado fica o corpo, no meio da floresta, no pé de uma casa de madeira no que aparenta ser um vilarejo no meio da floresta.